\documentclass[11pt,a4paper]{article}

\usepackage[margin=1in]{geometry}
\usepackage{amsmath,amssymb,bm}
\usepackage{booktabs}
\usepackage{array}
\usepackage{hyperref}
\usepackage{xcolor}
\usepackage{enumitem}

\hypersetup{
  colorlinks=true,
  linkcolor=black,
  urlcolor=blue!50!black,
  citecolor=black
}

\pagestyle{plain}

\title{\textbf{BeamLock}\\[0.4em]
\large Physics-Informed Virtual Coarse-to-Fine PAT\\
for Mobile Free-Space Optical Communication\\[0.6em]
\normalsize Technical Deep Dive, Results, and Oral Briefing Notes}
\author{Team Cadets\\
Smart India Hackathon 2026 · Problem ID SIH26169\\
Theme: Smart Automation · Category: Software\\
Repository: \texttt{https://github.com/MudithDShetty/SIH}}
\date{\today}

\begin{document}
\maketitle
\begin{abstract}
BeamLock (repository name: \emph{FSOC Virtual Tracker}) is a real-time software sandbox for Pointing, Acquisition, and Tracking (PAT) of mobile free-space optical (FSOC) links.
It closes a coarse AZ/EL camera loop under a statistical atmospheric channel (Cn$^2$ $\rightarrow$ Fried $r_0$, beam wander, scintillation), Gaussian-beam link SNR, Classical or AI+fusion beacon detection, constant-velocity Kalman tracking, raster re-acquisition when lock is lost, and a readiness-gated mock fine stage (4-quadrant detector + FSM).
This document explains the full system architecture, governing equations as implemented in code, control logic, evaluation methodology, verified numerical results, honest limitations, and a structured 20-minute briefing script.
It is intended for SIH evaluators and technical reviewers.
\end{abstract}

\tableofcontents
\newpage

%% ------------------------------------------------------------------
\section{Problem Statement and Scope}
%% ------------------------------------------------------------------

\subsection{SIH problem}
\textbf{ID:} SIH26169.\\
\textbf{Title:} Development of an AI-Based Virtual Camera Tracking System for Coarse Alignment of Mobile Free Space Optical Communication (FSOC) Terminals.\\
\textbf{Theme:} Smart Automation · \textbf{Category:} Software · \textbf{Team:} Cadets.

\subsection{Operational need}
FSOC links use narrow laser beams.
Before a fine pointing stage (fast steering mirror / FSM and quadrant detector / 4-QD) can close the communications beam, a \emph{coarse} stage must:
\begin{enumerate}[leftmargin=1.4em]
  \item observe a remote beacon with a wide-field camera;
  \item detect the beacon under turbulence, vibration, and sensor noise;
  \item estimate its motion;
  \item continuously slew the camera/gimbal so the beacon remains in the field of view (FOV).
\end{enumerate}
Hardware PAT benches are costly and slow for algorithm iteration.
BeamLock provides a laptop-scale virtual camera and closed-loop control environment with Classical vs.\ AI comparison and reproducible logs.

\subsection{What BeamLock is and is not}
\begin{center}
\begin{tabular}{@{}p{0.46\textwidth}p{0.46\textwidth}@{}}
\toprule
\textbf{Is} & \textbf{Is not} \\
\midrule
Real-time coarse PAT simulator ($\sim$60\,fps) & Flight-qualified ISRO/SDA software \\
Physics-informed statistical channel + SNR handoff score & Certified BER / outage calculator \\
Classical and AI+fusion detectors in-loop & Split-step / phase-screen wave-optics propagator \\
Raster re-acquisition + Fine 4-QD mock & Optical hardware-in-the-loop \\
CSV/JSON + Streamlit PDF evidence & A full communications modem \\
\bottomrule
\end{tabular}
\end{center}

%% ------------------------------------------------------------------
\section{System Architecture}
%% ------------------------------------------------------------------

\subsection{End-to-end data flow}
\begin{enumerate}[leftmargin=1.4em]
  \item \textbf{UI / presets:} Calm, UAV, Stress; sliders for turbulence, vibration, noise, slew; detector toggle; multi-beacon.
  \item \textbf{Scene:} one or more moving beacons (linear / circular / random walk).
  \item \textbf{Disturbances + physics:} Cn$^2$ channel, platform vibration AR(1), sensor noise; Gaussian link budget.
  \item \textbf{VirtualCamera:} FOV crop, AZ/EL plant with rate/acceleration limits, true IFOV.
  \item \textbf{Detection:} Classical \emph{or} AI+fusion (mutually exclusive live path).
  \item \textbf{Kalman tracker:} state LOCKED / COASTING / LOST; camera driven by estimate, not ground truth.
  \item \textbf{Link Readiness:} physics-informed score in $[0,1]$; LOST forces 0.
  \item \textbf{Control:}
  \begin{itemize}
    \item if LOST $\rightarrow$ raster re-acquisition waypoints + distance-scaled slew;
    \item else if Fine active $\rightarrow$ mock 4-QD / FSM residual nulling;
    \item else $\rightarrow$ coarse AZ/EL slew toward Kalman estimate.
  \end{itemize}
  \item \textbf{Outputs:} µrad-primary HUD, CSV + run-sheet JSON, Streamlit report/PDF, offline Monte Carlo / Pd--Pfa / ablation scripts.
\end{enumerate}

\subsection{Primary modules}
\begin{center}
\begin{tabular}{@{}ll@{}}
\toprule
\textbf{Module} & \textbf{Role} \\
\midrule
\texttt{scene/} & Beacon motion, virtual world render \\
\texttt{scene/camera.py} & AZ/EL plant, FOV, IFOV \\
\texttt{disturbance/} & Warp residual, vibration, sensor noise coupling \\
\texttt{physics/atmosphere.py} & Cn$^2\rightarrow r_0$, Rytov, wander, scintillation \\
\texttt{physics/link\_budget.py} & $P_{rx}$, SNR, fade margin, pointing loss \\
\texttt{detector/} & Classical, YOLO, Hybrid fusion, preprocess \\
\texttt{tracker/kalman\_tracker.py} & Constant-velocity KF + state machine \\
\texttt{control/reacquisition.py} & Raster (default) search \\
\texttt{control/fine\_pointing.py} & Readiness-gated Fine mock \\
\texttt{metrics/} & Link readiness, CSV/JSON logging \\
\texttt{ui/} & Controls and scenario presets \\
\texttt{report.py} & Streamlit analysis + PDF \\
\texttt{main.py} & Live closed loop \\
\bottomrule
\end{tabular}
\end{center}

\subsection{Scenario presets}
\begin{center}
\begin{tabular}{@{}lrrrll@{}}
\toprule
\textbf{Name} & \textbf{Turb} & \textbf{Vib} & \textbf{Noise} & \textbf{Slew} & \textbf{Trajectory} \\
\midrule
Calm (F1) & 0 & 0 & 0.05 & 90 & Circular \\
UAV (F2) & 4 & 8 & 0.2 & 90 & Circular \\
Stress (F3) & 7 & 15 & 0.35 & 120 & Random walk \\
\bottomrule
\end{tabular}
\end{center}

Demo keys: \texttt{T} Classical$\leftrightarrow$AI+fusion, \texttt{L} force detector blind, \texttt{B} 1$\leftrightarrow$2 beacons.

%% ------------------------------------------------------------------
\section{Angular Geometry and Camera Plant}
%% ------------------------------------------------------------------

\subsection{Instantaneous field of view (IFOV)}
Default acquisition FOV is scaled to the ESA IZN-1 class $\sim$2.5\,mrad horizontal FOV:
\begin{equation}
\mathrm{IFOV}
= \frac{\mathrm{FOV}_{H}}{N_{x}}
= \frac{2500\,\mu\mathrm{rad}}{400\,\mathrm{px}}
= 6.25\,\mu\mathrm{rad/px}.
\end{equation}
Pixel error maps to angular error by
\begin{equation}
\theta_{\mu\mathrm{rad}} = e_{\mathrm{px}} \cdot \mathrm{IFOV}.
\end{equation}

\subsection{AZ/EL plant}
The virtual camera maintains boresight $(x,y)$ in scene coordinates with rate/acceleration limits (slew rate exposed in UI as px/s and equivalently µrad/s via IFOV).
When tracking is healthy, \texttt{point\_towards} drives toward the Kalman estimate.
During LOST search, slew multipliers increase with distance to the current waypoint (capped by \texttt{MAX\_ACQ\_SLEW\_MULTIPLIER} in \texttt{main.py}).

%% ------------------------------------------------------------------
\section{Atmospheric Channel Physics}
%% ------------------------------------------------------------------

Implementation: \texttt{physics/atmosphere.py}.
References: Andrews \& Phillips style order-of-magnitude models for horizontal / UAV-class paths.
\textbf{Not} a split-step phase-screen propagator.

\subsection{Cn$^2$ mapping from UI strength}
Slider strength $s\in[0,10]$ maps log-uniform onto $[C_{n,\min}^2, C_{n,\max}^2]$ with
$C_{n,\min}^2=10^{-17}$, $C_{n,\max}^2=10^{-12}$\,m$^{-2/3}$:
\begin{equation}
C_n^2(s)
=
\begin{cases}
C_{n,\min}^2, & s\le 0,\\[0.3em]
10^{\,\log_{10} C_{n,\min}^2
+ (\log_{10} C_{n,\max}^2-\log_{10} C_{n,\min}^2)\,(s/10)}, & s>0.
\end{cases}
\end{equation}
Interpretation: calm $\sim 10^{-17}$, moderate $\sim 10^{-14}$, strong $\sim 10^{-12}$.

\subsection{Default optical / geometric assumptions}
\begin{center}
\begin{tabular}{@{}ll@{}}
\toprule
Symbol / quantity & Default \\
\midrule
Wavelength $\lambda$ & $976\,\mathrm{nm}$ (acquisition-beacon class) \\
Path length $L$ & $8000\,\mathrm{m}$ (UAV / short terrestrial) \\
Receive aperture $D$ & $0.08\,\mathrm{m}$ \\
\bottomrule
\end{tabular}
\end{center}

\subsection{Fried parameter $r_0$}
Plane-wave form used in code:
\begin{equation}
r_0
=
\bigl(0.423\, k^2 C_n^2 L\bigr)^{-3/5},
\qquad
k=\frac{2\pi}{\lambda}.
\end{equation}
Smaller $r_0$ implies worse atmospheric coherence / seeing.

\subsection{Rytov variance}
\begin{equation}
\sigma_R^2
=
1.23\, C_n^2\, k^{7/6}\, L^{11/6}.
\end{equation}
Used as the driver for scintillation sampling regime.

\subsection{Beam wander / tip--tilt RMS}
One-axis RMS wander (µrad) blends an angle-of-arrival scale $\sim 0.5\lambda/r_0$ with mild aperture averaging and a Cn$^2$-dependent floor so mid-strength turbulence remains visible:
\begin{align}
\sigma_{\mathrm{AoA}}
&\approx
0.5\,\frac{\lambda}{r_0}\cdot A(D,r_0),\\
\sigma_{\mathrm{floor}}
&\propto
\sqrt{C_n^2}\,\sqrt{L},\\
\sigma_{\mathrm{wander}}
&=
\max(\sigma_{\mathrm{AoA}},\sigma_{\mathrm{floor}}).
\end{align}
Temporal correlation uses an AR(1) update with coefficient $\alpha\approx 0.85^{\,60\Delta t}$ (Greenwood-ish tens-of-ms memory).

\subsection{Scintillation irradiance factor $I$}
Irradiance multiplier with mean $\approx 1$:
\begin{itemize}[leftmargin=1.4em]
  \item Weak ($\sigma_R^2<0.5$): lognormal sampling with $\sigma_I^2\approx\sigma_R^2$.
  \item Moderate/strong: gamma--gamma style product of two unit-mean gammas (simplified Andrews $\alpha,\beta$ mapping), clamped to a practical dynamic range for the 60\,fps loop.
\end{itemize}
Scintillation is temporally smoothed more slowly than tip/tilt (aperture-averaging intuition).

\subsection{Appearance coupling}
Wander offsets the drawn beacon centroid in scene pixels.
Scintillation and pointing loss modulate spot intensity.
Turbulence and sensor-noise models further degrade the camera-view image before detection.

%% ------------------------------------------------------------------
\section{Optical Link Budget and Link Readiness}
%% ------------------------------------------------------------------

\subsection{Received power model}
Implementation: \texttt{physics/link\_budget.py}.
Directed FSO form (cf.\ IISc/QOSMIC framing, arXiv:2507.20908):
\begin{equation}
P_{rx}
=
P_{tx}\,
\eta_{tx}\,
\eta_{rx}\,
G_{\mathrm{geo}}\,
T_{\mathrm{atm}}\,
L_{pe}\,
I_{\mathrm{scint}}.
\end{equation}
Default parameters (simulation, not a flight terminal datasheet):
$P_{tx}=0.5\,\mathrm{W}$,
$\eta_{tx}=\eta_{rx}=0.7$,
$D_{rx}=0.08\,\mathrm{m}$,
divergence half-angle $\theta_{\mathrm{div}}\approx 2.6\,\mathrm{mrad}$,
$L=8\,\mathrm{km}$,
static $T_{\mathrm{atm}}=0.75$,
NEP-like noise floor $5\times 10^{-9}\,\mathrm{W}$.

\subsection{Geometric collection}
\begin{equation}
G_{\mathrm{geo}}
=
\min\!\left(1,\frac{\pi (D_{rx}/2)^2}{\pi (\theta_{\mathrm{div}} L)^2}\right).
\end{equation}

\subsection{Gaussian pointing loss}
With residual angle $\theta$ (rad):
\begin{equation}
L_{pe}
=
\exp\!\left(-2\left(\frac{\theta}{\theta_{\mathrm{div}}}\right)^2\right).
\end{equation}

\subsection{SNR and fade margin}
\begin{align}
\mathrm{SNR}_{\mathrm{lin}}
&=
\frac{P_{rx}}{P_{\mathrm{noise}}},\\
\mathrm{SNR}_{\mathrm{dB}}
&=
10\log_{10}\mathrm{SNR}_{\mathrm{lin}},\\
M_{\mathrm{fade}}
&=
\mathrm{SNR}_{\mathrm{dB}}-10\,\mathrm{dB}.
\end{align}
The $10\,\mathrm{dB}$ term is a handoff floor used for readiness, \textbf{not} a claimed BER threshold.

\subsection{Link Readiness score}
Implementation: \texttt{metrics/link\_readiness.py}.
Physics path (default in \texttt{main.py}):
\begin{equation}
S
=
\mathrm{clamp}\!\big(
0.45\,s_{\mathrm{SNR}}
+
0.25\,L_{pe}
+
0.15\,\min(1,I)
+
0.15\,s_{\mathrm{lock}}
,\,0,\,1\big),
\end{equation}
with hard rule
\begin{equation}
\text{if tracking state $=$ LOST then } S=0.
\end{equation}
$s_{\mathrm{SNR}}$ maps roughly $10\,\mathrm{dB}\rightarrow 0$, $25\,\mathrm{dB}\rightarrow 1$.
\textbf{Purpose:} coarse-to-fine handoff decision support.
\textbf{Not:} certified BER, CCSDS outage, or flight link closure certificate.

%% ------------------------------------------------------------------
\section{Detection, Tracking, and Re-acquisition}
%% ------------------------------------------------------------------

\subsection{Classical detector}
Grayscale threshold (default 200) $\rightarrow$ largest contour centroid.
Baseline for a bright beacon on dark sky.
Adaptive mode available under elevated sensor noise.

\subsection{AI detector}
YOLOv8n trained on synthetic frames generated from the same disturbance pipeline
(\texttt{scripts/generate\_training\_data.py}, \texttt{scripts/train.py}).
Weights: \texttt{weights/beacon\_yolov8n.pt}.
Returns highest-confidence box centroid.

\subsection{AI+fusion (live ``AI'' path)}
\texttt{HybridDetector}:
\begin{enumerate}[leftmargin=1.4em]
  \item denoise / top-hat enhance;
  \item YOLO on raw + enhanced;
  \item adaptive classical on enhanced;
  \item fuse by consensus ($\sim$35\,px) or best candidate;
  \item short temporal hold (8 frames) during gaps.
\end{enumerate}
Raw YOLO alone is fragile in heavy grain; fusion is the production live path when AI is selected.

\subsection{Kalman tracker}
State $\mathbf{x}=[x,y,v_x,v_y]^\top$ with constant-velocity transition
\begin{equation}
F
=
\begin{bmatrix}
1 & 0 & \Delta t & 0\\
0 & 1 & 0 & \Delta t\\
0 & 0 & 1 & 0\\
0 & 0 & 0 & 1
\end{bmatrix},
\quad
H
=
\begin{bmatrix}
1 & 0 & 0 & 0\\
0 & 1 & 0 & 0
\end{bmatrix}.
\end{equation}
Process noise $Q$ uses the standard discrete white-noise acceleration structure scaled by process-noise intensity.
Measurement noise $R=\sigma_z^2 I_2$.

\subsection{Tracking state machine}
\begin{center}
\begin{tabular}{@{}ll@{}}
\toprule
\textbf{State} & \textbf{Condition} \\
\midrule
LOCKED & Valid accepted detection this frame \\
COASTING & 1--15 consecutive missed measurements \\
LOST & $>15$ misses (or forced lost on detector switch) \\
\bottomrule
\end{tabular}
\end{center}
Camera control uses the Kalman estimate when not LOST.
On LOST, re-acquisition owns the slew target.

\subsection{Re-acquisition}
\texttt{ReacquisitionController}:
\begin{itemize}[leftmargin=1.4em]
  \item \textbf{Default mode on activate: raster} (covariance-scaled step spacing).
  \item Spiral waypoint generator exists in code but is not the default activate path (spiral corner-trapped under high error).
  \item Search center biased by last lock + last velocity + lost time.
  \item Plausibility gating rejects false re-locks far from search/boresight while LOST.
\end{itemize}

%% ------------------------------------------------------------------
\section{Fine Pointing Stage}
%% ------------------------------------------------------------------

Implementation: \texttt{control/fine\_pointing.py}.

\subsection{Handoff criteria}
\begin{center}
\begin{tabular}{@{}ll@{}}
\toprule
Parameter & Value \\
\midrule
Enter readiness & $S\ge 0.70$ \\
Exit readiness & $S\le 0.55$ (hysteresis) \\
Tracking & must be LOCKED to enter \\
Capture cone enter & residual $\le 600\,\mu\mathrm{rad}$ \\
Capture drop & residual $>900\,\mu\mathrm{rad}$ \\
Stable frames & 12 ($\sim$0.2\,s at 60\,fps) \\
Fine slew authority & $4000\,\mu\mathrm{rad/s}$ class \\
Quad noise / FSM latency & mild demo non-idealities \\
\bottomrule
\end{tabular}
\end{center}

\subsection{4-QD mock}
Spot location relative to FOV center yields normalized quadrant imbalance $(q_x,q_y)$ and residual angle.
The controller nulls residual with high-bandwidth tip/tilt relative to the coarse plant.
This is a \textbf{gated demo plant}, not a calibrated flight FSM/4-QD identification.

%% ------------------------------------------------------------------
\section{Evaluation Methodology}
%% ------------------------------------------------------------------

\subsection{Matched-seed Classical vs AI}
\texttt{scripts/run\_comparison.py} reseeds per-frame RNG so turbulence/vibration/noise realizations match across detectors.
Primary metrics: lock retention \%, re-acquisition count, mean/max pixel error, readiness, acquisition time.

\subsection{Acquisition Monte Carlo}
\texttt{scripts/monte\_carlo\_acquisition.py} vs ESA ESTOL warm-start bound $T_{\mathrm{acq}}<60\,\mathrm{s}$.
Reports median, mean, p95, pass count.

\subsection{Clutter Pd/Pfa}
\texttt{scripts/evaluate\_detection.py} under Stress-like conditions.
Cite Classical numbers only when AI/hybrid JSON shows pathological $P_{fa}=1$.

\subsection{Ablation / packaging}
\texttt{scripts/run\_ablation.py}, Streamlit PDF from \texttt{report.py}, \texttt{scripts/build\_exe.py} for standalone packaging.
Demo video: \texttt{scripts/record\_demo\_video.py} $\rightarrow$ \texttt{logs/demo/beamlock\_sih\_demo.mp4}.

%% ------------------------------------------------------------------
\section{Verified Results}
%% ------------------------------------------------------------------

\subsection{Classical vs AI scenario matrix}
Source: \texttt{docs/technical\_report.md} (matched seed $\approx$45\,s, seed 42, after hard-curriculum retrain).

\begin{center}
\begin{tabular}{@{}lrrrrr@{}}
\toprule
Scenario & C err & AI err & C lock & AI lock & Reacq C/AI \\
\midrule
Calm & 35.5\,px & 35.5\,px & 100\% & 100\% & 1 / 1 \\
UAV & 35.6\,px & 35.6\,px & 97.7\% & \textbf{100\%} & 1 / 1 \\
Stress & \textbf{35.1}\,px & 36.1\,px & 95.5\% & 95.5\% & 2 / \textbf{1} \\
\bottomrule
\end{tabular}
\end{center}

\paragraph{Interpretation.}
Mean pixel error remains $\sim$35\,px even in Calm: largely \textbf{slew-lag limited}.
AI's documented advantage is lock retention (UAV) and fewer re-acquisitions (Stress), not mean-error domination.
Do not claim Stress mean-error superiority.

\subsection{Acquisition Monte Carlo (file-backed)}
\begin{center}
\begin{tabular}{@{}lrrrrr@{}}
\toprule
Scenario & Trials & Acquired & ESTOL$<$60\,s & Median $T_{\mathrm{acq}}$ & p95 \\
\midrule
UAV & 40 & 40/40 & 40/40 & 0.817\,s & 0.85\,s \\
Stress & 40 & 40/40 & 40/40 & 0.967\,s & 1.0\,s \\
\bottomrule
\end{tabular}
\end{center}
Files: \texttt{logs/monte\_carlo/acquisition\_uav\_40trials.json},
\texttt{logs/monte\_carlo/acquisition\_stress\_40trials.json}.

\subsection{Classical Stress clutter}
From \texttt{logs/detection/pd\_pfa\_stress.json}:
$P_d\approx 0.667$, $P_{fa}\approx 0.024$.
Do \textbf{not} cite AI/hybrid rows with $P_{fa}=1.0$ from that evaluation dump.

\subsection{Numbers deliberately excluded}
Chat-only fusion lock percentages not saved to logs (e.g.\ unverified 81.8\%$\rightarrow$100\% / 79.5\%$\rightarrow$97.7\%) must not appear in submissions until re-run and archived.

%% ------------------------------------------------------------------
\section{Standards Alignment and References}
%% ------------------------------------------------------------------

\begin{enumerate}[leftmargin=1.4em]
  \item SDA Optical Communications Terminal (OCT) Standard v3.0.1 --- coarse/fine PAT stages.
  \item ESA ESTOL air interface v2.2 (REQ-PHY-030) --- $T_{\mathrm{acq}}<60\,\mathrm{s}$ warm start.
  \item ESA IZN-1 --- $\sim$2.5\,mrad acquisition FOV class; µrad-class fine tracking class.
  \item Andrews \& Phillips, \emph{Laser Beam Propagation through Random Media} --- Cn$^2$, $r_0$, scintillation, wander.
  \item Kysel\'{a}k et al., \emph{Photonics} 11(6):540 (2024) --- FSO acquisition scan patterns (doi:10.3390/photonics11060540). Note: literature emphasizes spirals; BeamLock default is raster.
  \item Kohantorabi \& Amirabadi, \emph{Applied Optics} --- deep vision YOLO FSO-PAT (doi:10.1364/ao.595667).
  \item Shivkant, Jain, Ramakrishnan, arXiv:2507.20908 (IISc/QOSMIC) --- probabilistic optical link-budget framing.
  \item ISRO Opto-Quantum Communication (public TEC) --- architecture inspiration only.
\end{enumerate}

%% ------------------------------------------------------------------
\section{Limitations and Honest Claims}
%% ------------------------------------------------------------------

\begin{itemize}[leftmargin=1.4em]
  \item Statistical channel, not phase-screen propagation.
  \item Link Readiness $\neq$ BER.
  \item Fine 4-QD/FSM is a readiness-gated mock with mild noise/latency.
  \item Synthetic YOLO training distribution; on-sky transfer unproven here.
  \item Mean pixel error is a weak detector discriminator under slew limits.
  \item Netlify hosts a static judge site + browser teaser; full stack is desktop \texttt{python main.py}.
  \item Some older docs still say ``spiral''; code default is raster.
\end{itemize}

%% ------------------------------------------------------------------
\section{20-Minute Briefing Script (Detailed)}
%% ------------------------------------------------------------------

\subsection{Timing map}
\begin{center}
\begin{tabular}{@{}llp{0.48\textwidth}@{}}
\toprule
Clock & Block & Focus \\
\midrule
0:00--1:30 & Hook & SIH26169, why FSOC needs PAT \\
1:30--5:00 & Solution & Pipeline, UVP, architecture \\
5:00--9:00 & Technical + physics & Channel, SNR, control law \\
9:00--11:30 & Feasibility & Laptop ops, risks/mitigations \\
11:30--14:00 & Impact & Who benefits, comparison table \\
14:00--16:00 & Results & Matrix, Monte Carlo, Pd/Pfa \\
16:00--18:30 & Demo & Video or live Calm/UAV/Stress/L/T \\
18:30--20:00 & Close & References, honest limit, Q invite \\
\bottomrule
\end{tabular}
\end{center}

\subsection{Spoken script (expanded)}
\paragraph{0:00--1:30.}
``We are Team Cadets. BeamLock addresses SIH26169: AI-based virtual camera tracking for coarse alignment of mobile FSOC terminals.
Free-space optical links offer enormous bandwidth, but the beam is narrow.
On a moving platform, turbulence, vibration, and noise break lock.
Hardware PAT benches make algorithm iteration slow and expensive.
BeamLock is a laptop simulator of the coarse PAT loop with optional Fine handoff, Classical versus AI+fusion, and reproducible metrics.''

\paragraph{1:30--5:00.}
``Solution pipeline: scene and beacon; Cn$^2$ physics channel; AZ/EL camera with true IFOV 6.25 microradians per pixel; Classical or AI+fusion detection; Kalman LOCKED/COASTING/LOST; Link Readiness; then control---raster re-acq if lost, Fine 4-QD if handed off, else coarse slew.
Uniqueness: closed-loop AZ/EL sandbox, AI inside the live loop, readiness-gated Fine, ESTOL-aligned Monte Carlo.
Mapped to the PS: virtual camera, coarse alignment under mobile disturbances, AI assistance, and evidence logs.''

\paragraph{5:00--9:00.}
``Physics: UI strength maps to Cn$^2$, then Fried $r_0$, Rytov variance, AR beam wander, and scintillation.
Link budget forms $P_{rx}$ with geometric collection and Gaussian pointing loss, yielding SNR and fade margin.
Readiness blends SNR, pointing loss, scintillation, and lock state; LOST forces zero.
Control law as above.
Scenarios Calm, UAV, Stress.
This is Andrews-style statistics for real-time use, not a phase screen.''

\paragraph{9:00--11:30.}
``Runs on a laptop; Classical CPU fallback; CUDA optional.
Ops: main.py, F1/F2/F3, T, L, B, Streamlit report.
Risks: AI noise fragility mitigated by fusion; search corner-trap mitigated by raster; mean-pixel vanity metric mitigated by lock/re-acq/Tacq reporting; flight-hardware misunderstanding mitigated by honest sandbox framing.''

\paragraph{11:30--14:00.}
``Impact: faster PAT algorithm R\&D, education in µrad/Cn$^2$/$r_0$/Tacq vocabulary, reproducible evidence culture.
Comparison versus typical early sims: AZ/EL plant, fusion A/B, physics channel, SNR readiness, Fine gate, raster, multi-beacon, Monte Carlo.''

\paragraph{14:00--16:00.}
``Results: UAV lock 97.7\% Classical vs 100\% AI; Stress re-acq 2 vs 1; mean error $\sim$35\,px slew-limited.
Monte Carlo 40/40 within ESTOL 60\,s; median Tacq 0.82\,s UAV, 0.97\,s Stress.
Classical Stress $P_d\approx0.67$, $P_{fa}\approx0.024$.''

\paragraph{16:00--18:30.}
Play \texttt{beamlock\_sih\_demo.mp4} or live demo.
Narrate overview vs FOV, cyan detection, magenta Kalman, readiness, forced blind, raster recover, AI+fusion segment.

\paragraph{18:30--20:00.}
Cite SDA OCT, ESTOL, IZN-1, Andrews, Photonics 540, Applied Optics YOLO-PAT, arXiv:2507.20908.
Close: sandbox, not BER/flight plant.
Invite questions.

\subsection{Hard Q\&A crib}
\begin{itemize}[leftmargin=1.4em]
  \item \textbf{Flight-ready?} No---virtual PAT sandbox.
  \item \textbf{BER?} No---readiness is handoff support.
  \item \textbf{Why 35\,px?} Slew lag; use lock/re-acq/Tacq.
  \item \textbf{Spiral?} Literature yes; our default is raster.
  \item \textbf{Netlify full app?} No---desktop prototype; site is portal + teaser + video.
\end{itemize}

%% ------------------------------------------------------------------
\section{Reproducibility Commands}
%% ------------------------------------------------------------------

\begin{verbatim}
cd fsoc-virtual-tracker
python -m venv .venv
.venv\Scripts\activate          # Windows
pip install -r requirements.txt
python main.py
streamlit run report.py
python scripts/run_comparison.py
python scripts/monte_carlo_acquisition.py --trials 40 --scenario uav
python scripts/record_demo_video.py --fps 15 --copy-site
\end{verbatim}

%% ------------------------------------------------------------------
\section{Conclusion}
%% ------------------------------------------------------------------

BeamLock closes a standards-aware coarse PAT loop in software: statistical atmosphere, Gaussian link SNR, Classical/AI+fusion detection, Kalman tracking, raster re-acquisition, and readiness-gated Fine mock, with reproducible Monte Carlo and A/B evidence.
For SIH26169 it demonstrates an AI-assisted virtual camera tracking system for coarse alignment of mobile FSOC terminals, with explicit physics assumptions and explicit non-claims.

\vspace{1.5em}
\noindent\textit{End of document.}

\end{document}
